%% Section 7: conclusion and future work. Budget: about 0.4 page.

\section{Conclusion}
\label{sec:conclusion}

Denotational equality makes compiler proofs short as long as a pass
keeps the representation of values. With a graph-model semantics for
every language, nine of the ten passes of a compiler from a
call-by-value $\lambda$-calculus with recursive types to a typed
assembly language are short equational calculations, and we mechanized
the whole proof in Agda. When a pass changes the representation of
closures, equality gives way to a relation, and semantic reflection
needs a new one, the approximant logical relation, which relates a
finite target approximation to the whole source denotation. That pass
is where the denotational proof pays: it is several times the size of
the operational proofs of the same pass. In future work we plan to
report on a dynamically typed variant of the compiler and to repeat
the comparison in Rocq or Lean.
